% !PRINTED 234
\weberpage{234}
%<<CONT
\weq{10}{\begin{aligned}
\mu &\equiv \alpha = a\alpha' + c\beta'\\
\mu' &\equiv \beta = \partial\beta'
\end{aligned}\ (\mathrm{mod}\ n),}

where $\beta'$ denotes an arbitrary number coprime to $a$, and $\alpha'$ is so
determined that $\alpha$, $\beta$ become relatively prime; and it is also easy
to see [by §\,65, (3)] that, so long as $a$, $\partial$, $c$ remain the same,
the pairs of numbers $\mu$, $\mu'$ determined by (10) belong to the same
row.

We call the numbers $a$, $c$, $\partial$ (as in §\,27) the \gesperrt{transformation
numbers}, and $n$ the \gesperrt{degree of transformation}.

\gesperrt{The system of numbers $a$, $c$, $\partial$ is therefore completely
characteristic for a row}, and the number of rows is equal to the
number of these systems of numbers, from which there follows for the number
$\psi(n)$, (§\,63):

\weq{11}{\psi(n) = \Sigma \frac{a}{e}\varphi(e),}

when the sum is extended over all divisors $a$ of $n$.

It is of interest to prove this relation directly as well, in which case the
restriction to an odd $n$ can be dropped. If we now regard $\psi(n)$ as a
symbol for the sum (11), then, when $n'$, $n''$ are relatively prime, we
obtain first
\[
\psi(n')\psi(n'') = \psi(n'n''),
\]
and it thus only remains to determine the sum $\psi(n)$ for the case in which
$n=p^\pi$ is a power of a prime. In this case $e$ is equal to the smaller of
the two divisors $a$, $\partial$, and when $a=\partial$, $e=a$. If, then, we
separate off the first and last term of the sum $\psi(n)$, we obtain
\[
\psi(p^\pi) = 1 + p^\pi + \sum_{1<a\leqq\sqrt{n}}\varphi(a) + \sum_{\sqrt{n}<a<n}\frac{a}{\partial}\varphi(\partial).
\]
But we have
\[
\varphi(a) = a\frac{p-1}{p}, \qquad \varphi(\partial) = \partial\frac{p-1}{p},
\]
hence
\[
\begin{aligned}
\psi(p^\pi) &= 1 + p^\pi + \frac{p-1}{p}\sum_{1<a<n} a\\
&= 1 + p^\pi + (p-1)(1+p+p^2+\cdots p^{\pi-2}) = p^\pi\Bigl(1+\frac{1}{p}\Bigr),
\end{aligned}
\]
%CONT>>
